% Options for packages loaded elsewhere
\PassOptionsToPackage{unicode}{hyperref}
\PassOptionsToPackage{hyphens}{url}
\documentclass[
  ignorenonframetext,
]{beamer}
\newif\ifbibliography
\usepackage{pgfpages}
\setbeamertemplate{caption}[numbered]
\setbeamertemplate{caption label separator}{: }
\setbeamercolor{caption name}{fg=normal text.fg}
\beamertemplatenavigationsymbolsempty
% remove section numbering
\setbeamertemplate{part page}{
  \centering
  \begin{beamercolorbox}[sep=16pt,center]{part title}
    \usebeamerfont{part title}\insertpart\par
  \end{beamercolorbox}
}
\setbeamertemplate{section page}{
  \centering
  \begin{beamercolorbox}[sep=12pt,center]{section title}
    \usebeamerfont{section title}\insertsection\par
  \end{beamercolorbox}
}
\setbeamertemplate{subsection page}{
  \centering
  \begin{beamercolorbox}[sep=8pt,center]{subsection title}
    \usebeamerfont{subsection title}\insertsubsection\par
  \end{beamercolorbox}
}
% Prevent slide breaks in the middle of a paragraph
\widowpenalties 1 10000
\raggedbottom
\AtBeginPart{
  \frame{\partpage}
}
\AtBeginSection{
  \ifbibliography
  \else
    \frame{\sectionpage}
  \fi
}
\AtBeginSubsection{
  \frame{\subsectionpage}
}
\usepackage{iftex}
\ifPDFTeX
  \usepackage[T1]{fontenc}
  \usepackage[utf8]{inputenc}
  \usepackage{textcomp} % provide euro and other symbols
\else % if luatex or xetex
  \usepackage{unicode-math} % this also loads fontspec
  \defaultfontfeatures{Scale=MatchLowercase}
  \defaultfontfeatures[\rmfamily]{Ligatures=TeX,Scale=1}
\fi
\usepackage{lmodern}
\ifPDFTeX\else
  % xetex/luatex font selection
\fi
% Use upquote if available, for straight quotes in verbatim environments
\IfFileExists{upquote.sty}{\usepackage{upquote}}{}
\IfFileExists{microtype.sty}{% use microtype if available
  \usepackage[]{microtype}
  \UseMicrotypeSet[protrusion]{basicmath} % disable protrusion for tt fonts
}{}
\makeatletter
\@ifundefined{KOMAClassName}{% if non-KOMA class
  \IfFileExists{parskip.sty}{%
    \usepackage{parskip}
  }{% else
    \setlength{\parindent}{0pt}
    \setlength{\parskip}{6pt plus 2pt minus 1pt}}
}{% if KOMA class
  \KOMAoptions{parskip=half}}
\makeatother
\setlength{\emergencystretch}{3em} % prevent overfull lines
\providecommand{\tightlist}{%
  \setlength{\itemsep}{0pt}\setlength{\parskip}{0pt}}
% Slide layout defaults for warning-clean scientific decks.
\usepackage{etoolbox}
\IfFileExists{xurl.sty}{\usepackage{xurl}}{}
\IfFileExists{seqsplit.sty}{\usepackage{seqsplit}}{\newcommand{\seqsplit}[1]{#1}}
\protected\def\breakseq#1{\seqsplit{#1}}
\protected\def\breaktt#1{\begingroup\ttfamily\seqsplit{#1}\endgroup}
\setlength{\emergencystretch}{6em}
\tolerance=5000
\hbadness=10000
\hfuzz=1pt
\setlength{\tabcolsep}{2pt}
\AtBeginEnvironment{longtable}{\tiny\renewcommand{\arraystretch}{0.86}\setlength{\tabcolsep}{1pt}}
\AtBeginEnvironment{tabular}{\tiny\renewcommand{\arraystretch}{0.86}\setlength{\tabcolsep}{1pt}}

% Natbib and cross-reference fallbacks — slides don't load natbib
% or cleveref, but combined-PDF manuscript prose may emit these
% commands. The fallback renders citations as a bracketed key list
% and cross-references as detokenized labels so slides don't fail on
% undefined control sequences or raw underscores. \providecommand is
% a no-op if packages load later.
\providecommand{\citep}[1]{[#1]}
\providecommand{\citet}[1]{#1}
\providecommand{\citealp}[1]{#1}
\providecommand{\citeauthor}[1]{#1}
\providecommand{\citeyear}[1]{#1}
\providecommand{\cref}[1]{\texttt{\detokenize{#1}}}
\providecommand{\Cref}[1]{\texttt{\detokenize{#1}}}

\newtheorem{proposition}{Proposition}
\newtheorem{hypothesis}{Hypothesis}
\usepackage{bookmark}
\IfFileExists{xurl.sty}{\usepackage{xurl}}{} % add URL line breaks if available
\urlstyle{same}
% fallback for those not using the hyperref driver hyperxmp:
\makeatletter
\@ifundefined{xmpquote}{\newcommand{\xmpquote}[1]{#1}}{}
\makeatother
\hypersetup{
  hidelinks,
  pdfcreator={LaTeX via pandoc}}

\author{\texorpdfstring{}{}}
\date{}

\begin{document}

\section{Discussion}\label{sec:discussion}

\begin{frame}[fragile,allowframebreaks]{What the demonstration does and
does not show}
\protect\phantomsection\label{what-the-demonstration-does-and-does-not-show}
The demonstration shows that the registered analysis, run against seeded
data with a genuine effect (\texttt{Normal(0.8,\ 1)} treatment versus
\texttt{Normal(0,\ 1)} control), recovers a significant confirmatory
result: an observed mean difference of \texttt{1.003} with a two-sided
permutation p-value of \texttt{0.0005}. It does \textbf{not} show
anything about any real-world phenomenon. The data are synthetic and the
effect is injected by construction; the number exists to prove that the
locked plan, when executed, produces an auditable statistic --- not to
support a scientific claim. This is the honest reading of a template:
the value is a property of the code path, not evidence about the world.
\end{frame}

\begin{frame}[fragile,allowframebreaks]{Why the boundaries are enforced
in code}
\protect\phantomsection\label{why-the-boundaries-are-enforced-in-code}
Prose conventions for ``confirmatory versus exploratory'' are easy to
blur once results are in view. Here the boundary is a function contract:
\breaktt{compare\_analysis\_to\_registration} refuses unregistered
outcomes without a documented deviation,
\breaktt{build\_deviation\_ledger} assigns a severity to every executed
element, and \breaktt{build\_review\_packet} partitions confirmatory from
exploratory outcomes deterministically. The content hash on the frozen
registration makes silent edits to the locked plan detectable
(\texttt{hash\_mismatch}). Together these turn registered-report
discipline into checks that fail loudly rather than guidance that can be
quietly ignored.
\end{frame}

\begin{frame}[fragile,allowframebreaks]{Limitations and scope}
\protect\phantomsection\label{limitations-and-scope}
\begin{itemize}
\tightlist
\item
  The demonstration uses a single hypothesis and a single confirmatory
  outcome; real registrations will declare several, and the same helpers
  scale to them.
\item
  A permutation test is used because it is exact, deterministic, and
  dependency-free; forks may register any primary model by naming it in
  the plan.
\item
  Ethics-review and registered-report-stage metadata in the fixture are
  synthetic (\texttt{status:\ exempt},
  \breaktt{authority:\ synthetic-fixture-review}) and must be replaced
  with real values before any submission.
\end{itemize}
\end{frame}

\begin{frame}[fragile,allowframebreaks]{Using this template}
\protect\phantomsection\label{using-this-template}
Fork with \breaktt{scripts/audit/copy\_exemplar.py}, replace
\breaktt{data/example\_registration.json} with your own preregistration,
swap the synthetic generator in
\breaktt{src/registered\_report/demo\_study.py} for real data collection,
record stage and ethics metadata, and rerun the tests and figure
generation. Keep confirmatory claims tied to registered outcomes, record
every departure in the deviation ledger, and never let post-run result
prose rewrite the preregistered intent.
\end{frame}

\end{document}
